\chapter{Introduction}
\label{secIntro}

\linnet{} is an application for computing the transfer function of linear
electronic circuits. The computation is performed symbolically rather than
numerically, and the result is a formula rather than a number or a series
of numbers. The resulting formula is the Laplace transform describing the
dependencies of the voltages and currents in the circuit on the input
voltages and currents. 

% Locally define a figure representation specific to the table.
{
\newcommand{\incTabFig}[2]{\includegraphics[width=#2]{#1}}

% This define is related to the specifics of the array package; see
% https://texwelt.de/fragen/3401/zentrieren-text-in-tabelle (visited Sep
% 20, 2026) for more.
\newcolumntype{M}[1]{>{\centering\arraybackslash}m{#1}}

\setlength{\extrarowheight}{5pt} % Locally increase cells for device icons.
\begin{table}[t]
\begin{center}
\begin{tabular}{|l|c|M{2cm}|}

\hline

% The column headers.
%   \multicolumn is applied to overrule the column cell type, we have
% defined with M{2cm}. Without, the word "Symbol" is horizontally centered
% but vertically bottom alined and sits lower than the other column
% titles.
Device                               & ID          & \multicolumn{1}{c|}{Symbol}   \\
\hline
% This hline makes the column titles a separate rectangular box, comment out to get just a
% line as separator.
%\hline

% Table entries start here.
Resistor                             & \code{R}    & \incTabFig{DeviceR}{1.9cm}    \\ 
Conductance                          & \code{Y}    & \incTabFig{DeviceY}{1.9cm}    \\
Capacitor                            & \code{C}    & \incTabFig{DeviceC}{0.8cm}    \\ 
Inductivity                          & \code{L}    & \incTabFig{DeviceL}{2cm}      \\ 
Ideal operational amplifier (op-amp) & \code{OP}   & \incTabFig{DeviceOP}{1.3cm}   \\ 
Constant voltage source              & \code{U}    & \incTabFig{DeviceU0}{1cm}     \\
Voltage controlled voltage source    & \code{U(U)} & \incTabFig{DeviceU(u)}{1cm}   \\
Current controlled voltage source    & \code{U(I)} & \incTabFig{DeviceU(i)}{1cm}   \\
Constant current source              & \code{I}    & \incTabFig{DeviceI0}{0.4cm}   \\
Voltage controlled current source    & \code{I(U)} & \incTabFig{DeviceI(u)}{0.4cm} \\
Current controlled current source    & \code{I(I)} & \incTabFig{DeviceI(i)}{0.4cm} \\
Current probe (wire)                 & \code{PI}   & \incTabFig{DevicePI}{1.8cm}   \\ 
\hline
\end{tabular}
\caption{Supported linear devices}
\label{tabSupportedDevices}
\end{center}
\end{table}
} % End of local scope of \incTabFig and \extrarowheight

A linear electronic circuit is a combination of the supported basic
devices listed in table~\ref{tabSupportedDevices}. The circuit is input to
\linnet{}. The representation of the circuit is a list of devices with
connectivity information. The interconnections are expressed by references
to nodes, where a node is a point in the circuit to which normally at
least two devices are connected. This leads to a simple formal syntax, the
circuit network list, or simply netlist. This list can be created and
maintained with a text editor; there is no graphical interface for editing
a circuit.

\begin{figure}
\centering
\includegraphics[width=5.29cm]{FirstIntroExample}
\caption{Simple example of a linear electronic circuit}
\label{figFirstIntroExample}
\end{figure}

The computed formulas are printed to the console and to the application
log file and can be used for further investigation, publications, or
didactic purposes.

To make the application somewhat more attractive, it also exports the
computed formulas as Octave or MATLAB script code. Numeric evaluation then
becomes a simple one-line command in Octave. The formulas are exported as
LTI transfer function objects so that the complete set of analysis
functions from the Octave Control package can be applied directly. This
ranges from simple transfer function plotting to stability analysis and
the computation of system responses to arbitrary inputs.

Please refer to figure~\ref{figFirstIntroExample} for an example of how
\linnet{} works. This is a simple RLC circuit with a second-order transfer
function. It can be represented by the following circuit netlist:

\begin{verbatim}
U Uin in  gnd
L L   in  K1
C C   K1  out
R R   out gnd
PLOT G U_out U_in
\end{verbatim}

\noindent
If this netlist is stored in file \file{rlc.cnl}, we can run \linnet{}:

\begin{verbatim}
linNet -o rlc.cnl
\end{verbatim}

\noindent
This would yield the output:

\begin{verbatim}
User-defined result G (Bode plot):
The dependency of U_out on U_in:
  U_out(s) = N_U_out_U_in(s)/D_U_out_U_in(s) * U_in(s), with
    N_U_out_U_in(s) = R*C * s
    D_U_out_U_in(s) = L*C * s^2
                      +R*C * s
                      +1
\end{verbatim}

\begin{figure}
\centering
\includegraphics[width=0.7\textwidth]{FirstIntroExample_OctaveOutput}
\caption{Octave plot of transfer function of the system in figure
\ref{figFirstIntroExample}}
\label{figFirstIntroExample_OctaveOutput}
\end{figure}

\noindent
Going to Octave and typing \code{G} to plot the transfer function (still
using the default device values) gives us
figure~\ref{figFirstIntroExample_OctaveOutput}. Additional plots or plots
with modified device values are then a matter of single commands in
Octave.

\linnet{} means ``linear network''. It is based on a symbolic solver for
linear equation systems. There is no way to model non-linear effects such
as noise, voltage or current limits, non-linear distortions, or switching
operations. All of these effects play an important role in real electronic
circuits, and many circuits even use these effects as their principle of
operation. You will not find \linnet{} helpful for the investigation of
such effects or circuits. There are many numeric circuit simulation tools
that are capable of this, first and foremost the popular open-source tool
SPICE and its many derivatives. \linnet{} is conceptually not a competitor
to these tools, although it can behave somewhat similarly when using
Octave as a numeric post-processor. \linnet{} is not the worse SPICE;
\linnet{} is different. 